\section{Peano 算术}


\subsection{结构与理论要分开}


标准自然数结构写作

\[
\mathcal N=(\mathbb N;0,S,+,\cdot).
\]

说 $\mathcal N\models\mathrm{PA}$ 是说这个结构满足 PA 的所有公理；说 PA 是理论，则是说由一阶公理和推理规则得到的句子集合。结构不是理论，$\mathrm{PA}$ 也不等于 $\operatorname{Th}(\mathcal N)$。

\subsection{一组常用公理}


在语言 $(0,S,+,\cdot)$ 中，PA 包含：

\[
\begin{aligned}
&\forall x\,S(x)\ne0,\\
&\forall x\forall y\,(S(x)=S(y)\to x=y),\\
&\forall x\,(x+0=x),\qquad
\forall x\forall y\,(x+S(y)=S(x+y)),\\
&\forall x\,(x\cdot0=0),\qquad
\forall x\forall y\,(x\cdot S(y)=x\cdot y+x).
\end{aligned}
\]

最关键的是\textbf{一阶归纳公理模式}：对每个公式 $\varphi(x,\vec z)$，都有

\[
\bigl(\varphi(0,\vec z)\land\forall x(\varphi(x,\vec z)\to\varphi(S(x),\vec z))\bigr)
\to\forall x\varphi(x,\vec z).
\]

这里每个公式产生一个公理，不是一个能量化“所有公式”的单一一阶句子。

\subsection{做题时的区分}


\begin{itemize}
\item 集合论中的自然数归纳说某个集合是归纳集；
\item PA 的归纳模式说某个一阶公式在理论中可归纳；
\item 标准模型中的数学归纳法是语义事实。
\end{itemize}


三者结论相近，但层次不同，不能把其中一个直接当作另外两个的公理。
